%% gia-handout-doc.tex
%% Copyright 2026 André Hilbig and Frederick Heyder
%
% This work may be distributed and/or modified under the
% conditions of the LaTeX Project Public License, either version 1.3
% of this license or (at your option) any later version.
% The latest version of this license is in
%   http://www.latex-project.org/lppl.txt
% and version 1.3 or later is part of all distributions of LaTeX
% version 2005/12/01 or later.
%
% This work has the LPPL maintenance status `maintained'.
%
% The Current Maintainer of this work is André Hilbig.
%
% This work consists of the files
%   gia-handout.cls and
%   gia-handout-doc.tex.

% --- PDF tagging setup (currently disabled; pending listings compatibility) ---
% \DocumentMetadata{
    % tagging = on,
    % tagging-setup = {
        % activate/all,
        % tabsorder = structure
    % },
    % pdfversion = 2.0,
    % lang = de,
    % pdfstandard = a-4f,
    % pdfstandard = ua-2
% }
% --- end of PDF tagging setup ---

\documentclass[a4paper]{article}

\usepackage[utf8]{inputenc}
\usepackage[T1]{fontenc}

\usepackage[ngerman]{babel}

\usepackage[german=guillemets]{csquotes}

\usepackage[
    backend=biber,
    style=authoryear
]{biblatex}
\begin{filecontents*}{\jobname.bib}
@online{fischer2024,
    author       = {Ulrike Fischer},
    title        = {How to hide a QR-Code from tagpdf and provide alt-text},
    date         = {2024-05-19},
    url          = {https://tex.stackexchange.com/q/718239},
    urldate      = {2026-03-20},
    organization = {Stack Exchange}
}
@manual{tagpdf2026,
    author       = {Ulrike Fischer},
    title        = {The \texttt{tagpdf} package -- package documentation, v0.99y},
    date         = {2026-01-29},
    url          = {https://ctan.org/pkg/tagpdf},
    urldate      = {2026-03-20},
    organization = {CTAN}
}
\end{filecontents*}
\addbibresource{\jobname.bib}

\usepackage{microtype}

\usepackage{graphicx}

\usepackage{xcolor}

% --- gia-handout.cls extract: GI colors & GI commands ---
\definecolor{giaBlue}{RGB}{14,32,80}
\definecolor{giaYellow}{RGB}{255,230,0}
\definecolor{giaLightYellow}{RGB}{255,242,204}
\definecolor{giaGray}{RGB}{74,74,74}

\newcommand{\GI}{GI}
\newcommand{\gi}{Gesellschaft f\"ur Informatik e.\,V.}
% --- end of gia-handout.cls extract ---

% --- enumitem disabled (pending tagpdf compatibility) ---
% \usepackage{enumitem}
% \begin{description}[style=nextline, leftmargin=1em]
    % \item[]
% \end{description}
% --- end of enumitem disabled ---

\usepackage{multicol}

\usepackage{tikz}
\usetikzlibrary{calc, positioning}

\usepackage{tcolorbox}
\tcbuselibrary{breakable, skins}

\usepackage{listings}
\lstset{
    basicstyle=\ttfamily,
    columns=fullflexible,
    language=[LaTeX]TeX,
    escapechar=*,
    commentstyle=\color{green!50!black}\bfseries
}

\usepackage{hyperref}

\hypersetup{
    % pdftitle={The gia-handout package, v0.99a},
    % pdfsubject={Documentation of the gia-handout package},
    pdftitle={Das gia-handout-Paket, v0.99a},
    pdfsubject={Dokumentation des Pakets gia-handout},
    pdfkeywords={latex, gi, gia-handout},
    pdfauthor={Andr\'e Hilbig, Frederick Heyder}
}

\usepackage{url}

\newcommand{\mypdflatex}{\texttt{pdflatex}}
\newcommand{\mylualatex}{\texttt{lualatex}}
\newcommand\pkg[1]{\texttt{#1}}
\newcommand\myopt[2]{%
    \textcolor{giaYellow!75!giaGray}{\texttt{#1}}%
    \if\relax\detokenize{#2}\relax\else
        \textnormal{ = #2}
    \fi
}
\newcommand\mycmd[3]{%
    \textcolor{giaYellow!75!giaGray}{\texttt{\textbackslash#1}}%
    \if\relax\detokenize{#2}\relax\else
        \textnormal{\texttt{\{}#2\texttt{\}}}%
    \fi
    \if\relax\detokenize{#3}\relax\else
        \textnormal{\texttt{[}#3\texttt{]}}%
    \fi
}
\newcommand\myenv[3]{%
    \textcolor{giaYellow!75!giaGray}{\texttt{\textbackslash{begin}\{#1\}}}%
    \if\relax\detokenize{#2}\relax\else
        \textnormal{\texttt{\{}#2\texttt{\}}}%
    \fi
    \if\relax\detokenize{#3}\relax\else
        \textnormal{\texttt{[}#3\texttt{]}}%
    \fi
}

\author{Andr\'e \textsc{Hilbig} \qquad Frederick \textsc{Heyder}}
\title{\pkg{gia-handout}}
\date{20.~März~2026}

% --- gia-handout.cls extract: giacolorbox environment ---
\newlength{\giacolorboxskip}
\setlength{\giacolorboxskip}{4mm}
\addtolength{\giacolorboxskip}{3pt}
\NewDocumentEnvironment{giacolorbox}{o}
{
    \begin{tcolorbox}[%
        colback=white,
        colframe=giaBlue!0!giaYellow,
        sharp corners=all,
        title=#1,
        fonttitle=\color{black},
        enhanced jigsaw,
        grow sidewards by=2mm,
        borderline={4mm}{-4mm}{giaGray!5!white},
        borderline={0.5pt}{0pt}{giaYellow},
        before skip=\giacolorboxskip,
        after skip=\giacolorboxskip,
        \IfNoValueT{#1}{notitle}
    ]
}{
    \end{tcolorbox}
}
% --- end of gia-handout.cls extract ---

\begin{document}

\begin{centering}

    \vspace*{0.03\textheight}

    {\Large\textsc{Klassendokumentation}}\par
    \vspace{1.5\baselineskip}

    {\huge\bfseries\makeatletter\@title\makeatother}\par
    \vspace{1\baselineskip}

    {\Large\LaTeX-Klasse für Veröffentlichungen der \GI}\par
    \vspace{1.2\baselineskip}

    {\large Version~0.99a -- 4.~Mai~2026}\par
    \vspace{2\baselineskip}

    {\large\makeatletter\@author\makeatother}\par
    \vspace{1.5\baselineskip}

    {\normalsize\href{mailto:hilbig@uni-wuppertal.de}{\nolinkurl{hilbig@uni-wuppertal.de}}}\par
    \vspace{0.67\baselineskip}

    {\normalsize\url{https://git.uni-wuppertal.de/ddi/material/gia-handout}}

\end{centering}

\vfill

\begin{abstract}
Die Klasse bündelt die Bemühungen des FA~IBS\footnote{Fachausschusses \enquote{Informatische Bildung in Schulen} der \GI}, der FG~DDI\footnote{Fachgruppe \enquote{Didaktik der Informatik} der \GI} und insbesondere von Andr\'e Hilbig, Veröffentlichungen der \GI\ im selben Stil wie die Bildungsstandards zu vereinheitlichen.
Hierbei werden spezifische Designvorgaben und Layouts berücksichtigt, um den Anforderungen des Corporate Design der \GI\ zu entsprechen.
\end{abstract}

\vfill

\begin{multicols}{2}
    \tableofcontents
\end{multicols}

\clearpage

\section{Lizenz und Anforderungen}
\subsection{Lizenz}
Diese Klasse darf nach den Bedingungen der LaTeX Project Public Lizenz, entweder Version 1.3 oder (nach Ihrer Wahl) jede spätere Version, verteilt und/oder verändert werden.
Die neueste Version dieser Lizenz findet sich unter \url{https://www.latex-project.org/lppl/}.
\subsection{Anforderungen}
\subsubsection{Packages}
Die Klasse \pkg{gia-handout} verwendet die folgenden Pakete, wobei derzeit nicht alle vollständig mit \pkg{tagpdf} kompatibel sind.\footnote{Vgl.~\url{https://latex3.github.io/tagging-project/tagging-status/}}
\begin{table}[!h]
    \centering
    \begin{tabular}{ll}
        \hline
        \textbf{Package} & \textbf{Kompatibilität mit \pkg{tagpdf}}\\
        \hline
        iftex & kompatibel\\
        fontenc & kompatibel\\
        fontspec & kompatibel\\
        babel & nicht geprüft\\
        csquotes & teilweise kompatibel\\
        microtype & kompatibel\\
        setspace & kompatibel\\
        geometry & kompatibel\\
        scrlayer-scrpage & nicht geprüft\\
        xcolor & kompatibel\\
        tikz & teilweise kompatibel\\
        pgfkeys & kompatibel\\
        pgfopts & kompatibel\\
        tcolorbox & teilweise kompatibel\\
        xparse & kompatibel\\
        hyperref & teilweise kompatibel\\
        bookmark & kompatibel\\
        qrcode & derzeit inkompatibel\\
        FiraSans & kompatibel\\
        FiraMono & kompatibel\\
        \hline
    \end{tabular}
\end{table}
\subsubsection{Libraries}
Zusätzlich werden die nachstehenden Libraries benötigt.
\begin{multicols}{2}
    \paragraph{\pkg{tikz}-Libraries}\hfill%
    \begin{itemize}
        \item \pkg{calc}
        \item \pkg{positioning}
    \end{itemize}
    \columnbreak
    \paragraph{\pkg{tcolorbox}-Libraries}\hfill%
    \begin{itemize}
        \item \pkg{breakable}
        \item \pkg{skins}
    \end{itemize}
\end{multicols}

\addtocontents{toc}{\protect\setcounter{tocdepth}{1}}

\section{Minimalbeispiel}
Das folgende Minimalbeispiel veranschaulicht die grundlegende Verwendung, einschließlich der Definition der Autoren und des koordinierenden Autors.
\begin{lstlisting}
\documentclass[autorenListe,koordination]{gia-handout}

\GIAsetzeautorenListe{%
    \mbox{Max Mustermann (Berlin)}
    *\dots*
}
\GIAsetzekoordination{%
    Max Mustermann
}

\begin{document}
*\dots*
\end{document}
\end{lstlisting}

\section{Klassenoptionen}
Folgende Klassenoptionen stellt die Klasse \pkg{gia-handout} bereit:
\begin{description}
    \item[\myopt{autorenListe}{}] \hypertarget{opt:autorenListe}{}\phantom{}\hfill{Default: \enquote{keine Autoren}}\newline
        Ermöglicht, eine Autorenliste für das Dokument zu setzen. Die Zuweisung erfolgt über \hyperlink{cmd:GIAsetzeautorenListe}{\mycmd{GIAsetzeautorenListe}{}{}}, die Ausgabe über \hyperlink{cmd:GIAautorenListe}{\mycmd{GIAautorenListe}{}{}}.
    \item[\myopt{koordination}{}] \hypertarget{opt:koordination}{}\phantom{}\hfill{Default: \enquote{keine Angabe}}\newline
        Ermöglicht, den koordinierenden Autor zu setzen. Die Zuweisung erfolgt über \hyperlink{cmd:GIAsetzekoordination}{\mycmd{GIAsetzekoordination}{}{}}, die Ausgabe über \hyperlink{cmd:GIAkoordination}{\mycmd{GIAkoordination}{}{}}.
    \item[\myopt{frontcover}{$\langle$filename$\rangle$}] \hypertarget{opt:frontcover}{}\phantom{}\newline optional einzubindende Vorderseite.
    \item[\myopt{backcover}{$\langle$filename$\rangle$}] \hypertarget{opt:backcover}{}\phantom{}\newline optional einzubindende Rückseite.
\end{description}
Werden Vorder- und/oder Rückseite mit den entsprechenden Klassenoptionen nicht explizit eingebunden, werden Standardseiten gesetzt (s.~Abb.~\ref{fig:frontcover}~und~\ref{fig:backcover}, S.~\pageref{fig:frontcover}).
\begin{figure}[!h]
    \centering
    \begin{minipage}{.45\textwidth}
        \centering
        \fbox{\includegraphics[width=\textwidth]{gia-default-frontcover.pdf}}
        \caption{Default \protect\hyperlink{opt:frontcover}{\texttt{frontcover}}}
        \label{fig:frontcover}
    \end{minipage}
    \hfill
    \begin{minipage}{.45\textwidth}
        \centering
        \fbox{\includegraphics[width=\textwidth]{gia-default-backcover.pdf}}
        \caption{Default \protect\hyperlink{opt:backcover}{\texttt{backcover}}}
        \label{fig:backcover}
    \end{minipage}
\end{figure}

\section{PDF-Tagging}
Die Klasse unterstützt \pkg{tagpdf} für barrierefreies PDF-Tagging mit \mypdflatex\ und \mylualatex. Weitere Hinweise finden sich in der \pkg{tagpdf}-Dokumentation\footfullcite{tagpdf2026}.
\begin{minipage}{\linewidth} % prevent page break
Zur Aktivierung ist \texttt{\textbackslash DocumentMetadata} vor \texttt{\textbackslash documentclass} einzufügen:
\begin{lstlisting}
\DocumentMetadata{
    tagging-setup = {activate/all, tabsorder = structure},
    pdfversion = 2.0,
    pdfstandard = a-4f,
    pdfstandard = ua-2,
    lang = de
}
\end{lstlisting}
\end{minipage}
\subsection{Workaround für \pkg{qrcode}}
Da das Paket \pkg{qrcode} derzeit nicht mit \pkg{tagpdf} kompatibel ist, implementiert die Klasse \pkg{gia-handout} einen Workaround von Ulrike~Fischer\footfullcite{fischer2024}.

\section{Corporate Design}
Gemäß dem Corporate Design der \GI\ werden Fira~Sans als serifenlose Standard\-schriftart sowie Fira~Mono als Festbreitenschrift geladen.\\

\noindent Zudem werden die folgenden vier Farben definiert:
\begin{multicols}{2}
    \begin{itemize}
        \item[\fcolorbox{black}{giaBlue}{\phantom{\rule{1em}{.67em}}}] \texttt{giaBlue}
        \item[\fcolorbox{black}{giaYellow}{\phantom{\rule{1em}{.67em}}}] \texttt{giaYellow}
        \item[\fcolorbox{black}{giaLightYellow}{\phantom{\rule{1em}{.67em}}}] \texttt{giaLightYellow}
        \item[\fcolorbox{black}{giaGray}{\phantom{\rule{1em}{.67em}}}] \texttt{giaGray}
    \end{itemize}
\end{multicols}

\section{Verfügbare Befehle}
Die Klasse \pkg{gia-handout} definiert die folgenden Befehle:
\begin{description}
    \item[\mycmd{hervor}{$\langle$arg$\rangle$}{}]
        \phantom{}\newline setzt den Text fett und in der Farbe \texttt{giaBlue}.
    \item[\mycmd{GI}{}{}]
        \phantom{}\newline ein Wrapper für \verb|GI|
    \item[\mycmd{gi}{}{}]
        \phantom{}\newline ein Wrapper für \verb|Gesellschaft f\"ur Informatik e.\,V.|
    \item[\mycmd{GIAsetzeautorenListe}{$\langle$arg$\rangle$}{}]
        \hypertarget{cmd:GIAsetzeautorenListe}{}\phantom{}\newline definiert die zu setzenden Autoren.
    \item[\mycmd{GIAautorenListe}{}{}]
        \hypertarget{cmd:GIAautorenListe}{}\phantom{}\newline setzt die Autoren, sofern \hyperlink{opt:autorenListe}{\myopt{autorenListe}{}} aktiviert ist.
    \item[\mycmd{GIAsetzekoordination}{$\langle$arg$\rangle$}{}]
        \hypertarget{cmd:GIAsetzekoordination}{}\phantom{}\newline definiert die zu setzende Koordination.
    \item[\mycmd{GIAkoordination}{}{}]
        \hypertarget{cmd:GIAkoordination}{}\phantom{}\newline setzt die Koordination, sofern \hyperlink{opt:koordination}{\myopt{koordination}{}} aktiviert ist.
    \item[\mycmd{GIAmitwirkende}{}{}]
        \phantom{}\newline setzt die Mitwirkenden.
\end{description}

\section{Verfügbare Umgebungen}
Die Klasse \pkg{gia-handout} definiert die folgenden Umgebungen:
\subsection{Mehrspaltiger Text}
\begin{description}
    \item[\myenv{zweispaltig}{}{}]
        \phantom{}\newline setzt den Text zweispaltig.
\end{description}
\subsection{\texttt{giacolorbox}}
\begin{description}
    \item[\myenv{giacolorbox}{}{$\langle$title$\rangle$}]
        \phantom{}\newline setzt eine benutzerdefinierte \texttt{tcolorbox}, die als \texttt{Div} getaggt wird.
\end{description}
\begin{minipage}{.45\textwidth}
    \begin{lstlisting}
\begin{giacolorbox}
    body
\end{giacolorbox}
    \end{lstlisting}
\end{minipage}
\hfill
\begin{minipage}{.45\textwidth}
    \centering
    \begin{giacolorbox}
    body
    \end{giacolorbox}
\end{minipage}
\begin{minipage}{.45\textwidth}
    \begin{lstlisting}
\begin{giacolorbox}[title]
    body
\end{giacolorbox}
    \end{lstlisting}
\end{minipage}
\hfill
\begin{minipage}{.45\textwidth}
    \centering
    \begin{giacolorbox}[title]
    body
    \end{giacolorbox}
\end{minipage}

\section{Weiterentwicklung}
\begin{itemize}
    \item Barrierefreie Dokumentation
    \item Generalisierte Lizenzvorlagen
    \item Mehrseitiger Einband mit \pkg{pdfpages}
    \item Prüfung einer mehrsprachigen Ausrichtung der Klasse
\end{itemize}

\section{Changelog}
\begin{description}
    \item[Version~0.99 -- 20.~März~2026]
        \phantom{}\newline Pre-Release: Grundlegende Klassendefinition.
    \item[Version~0.99a -- 4.~Mai~2026]
        \phantom{}\newline Rename: Umbenennung von \pkg{gi-handout} zu \pkg{gia-handout}.
    \item[Version~1.00 -- tba]
        \phantom{}\newline \textit{Initial Release: Barrierefreie Dokumentation.}
\end{description}

\end{document}

\endinput